% ============================================================================
\clearpage
\section{Estado del proyecto al 28 de septiembre de 2026}\label{sec:estado}
% ============================================================================

Las secciones anteriores describen el sistema tal como se propuso el 26 de septiembre. Esta sección cuenta
qué se construyó desde entonces, qué funciona ya, qué cambió respecto de la propuesta y qué falta. Refleja
el estado de la mañana del lunes 28 de septiembre de 2026, antes de la primera sesión capturada de forma
programada. El detalle de cada paso está en los documentos del repositorio: el resumen de la sesión con
Alpaca, la descripción de la fuente y las respuestas a cada revisión.

\begin{ideaclave}[En pocas palabras]
\begin{itemize}
\item \textbf{La idea.} Cada mañana, a las 09:45 de Nueva York, se mide cuánto más cuesta protegerse de una
caída del S\&P~500 que apostar a una subida (el RR25 de las opciones SPXW). Después se comprueba si su
cambio de un día a otro anticipa la dirección o el riesgo de la bolsa, medida con SPY.
\item \textbf{Lo logrado.} Hay fuente de datos, un programa que mide y calcula lo que la bolsa hizo después,
y un proceso automático que guarda los precios cada día hábil. Está instalado, probado y confirmado por una
verificación independiente.
\item \textbf{Lo que falta.} Vigilar las primeras 3--5 sesiones, juntar 20--40 para el piloto y, con bastante
más historia, hacer la prueba predictiva con reglas fijadas de antemano.
\item \textbf{Lo que todavía no se sabe.} Si la señal predice algo. No hay ningún resultado sobre eso.
\end{itemize}
\end{ideaclave}

\subsection{Qué se construyó}\label{sec:estado_construido}

La \tabref{tab:construido} resume las piezas. Las comprueban 203 pruebas automáticas que no usan la red. La
batería también pasa en Windows y con Python obligado a declarar la codificación de cada texto.

\begin{table}[H]
\centering\small
\caption{Piezas construidas y su estado.}
\label{tab:construido}
\begin{tabularx}{\linewidth}{@{}>{\RaggedRight}p{3.3cm} Y >{\RaggedRight\arraybackslash}p{3.0cm}@{}}
\toprule
\encab{Pieza} & \encab{Qué hace} & \encab{Estado}\\
\midrule
Núcleo de referencia & Los métodos de este documento, con filtros causales, cuantiles con mesetas y
volatilidad implícita robusta. & Corregido\\
Pipeline del piloto & Ingesta, controles, medidas (RR25, asimetría y paridad), etiquetas, informe y
manifiesto reproducible. & Listo, con plantilla sintética\\
Adaptador de Alpaca & Reintentos, crudo inmutable con sus hashes, normalización al contrato y nivel
implícito de SPX. & Verificado con datos reales\\
Captura diaria & Ráfaga 5~s antes de las 09:45 y de las 10:00, página a página, con recuperación y un
estado por hora. & Instalada el 27 de septiembre\\
Histórico del objetivo & SIP de SPY en cada corte, 15 minutos después, y dividendos de SPY. & Instalado el
27 de septiembre\\
Etiquetas y dividendos & Rendimiento total y de precio de SPY, con dividendos versionados y sus
discrepancias. & Listo\\
Revisión de la operación & Estado, puntualidad, respaldo, SIP y dividendos de cada corte. & En uso desde el
28 de septiembre\\
\bottomrule
\end{tabularx}
\end{table}

\subsection{Datos: Alpaca y la captura hacia adelante}\label{sec:estado_datos}

Alpaca se verificó el 26 de septiembre con su API y su documentación (\tabref{tab:alpaca}). El hallazgo que
más condiciona el plan es que \textbf{no guarda cotizaciones pasadas de opciones}: el bid/ask de las 09:45
de sesiones anteriores no se puede reconstruir. Por eso la muestra se captura hacia adelante, desde el 28 de
septiembre.

\begin{table}[H]
\centering\small
\caption{Qué ofrece Alpaca para este proyecto (verificado el 26 de septiembre de 2026).}
\label{tab:alpaca}
\begin{tabularx}{\linewidth}{@{}>{\RaggedRight}p{3.6cm} Y@{}}
\toprule
\encab{Tema} & \encab{Resultado}\\
\midrule
SPXW & Disponible: opciones europeas, liquidación PM y vencimientos todos los días hábiles.\\
Historia de cotizaciones & \textbf{No hay.} Solo barras y operaciones de opciones desde febrero de 2024.\\
Nivel de SPX & \textbf{No hay.} Se estima por paridad con el vencimiento SPXW más cercano.\\
Feed gratuito (\texttt{indicative}) & Según Alpaca, sus cotizaciones están \emph{modificadas}: no son el NBBO
de OPRA.\\
Feed OPRA & Requiere el plan Algo Trader Plus (99~USD al mes) y firmar el acuerdo de OPRA.\\
SPY & Cotización IEX en tiempo real; NBBO consolidado (SIP) consultable 15 minutos después; dividendos en los
eventos corporativos, sin garantía de cuándo se publican.\\
\bottomrule
\end{tabularx}
\end{table}

La primera verificación con datos reales usó las últimas cotizaciones del viernes 25 al cierre. No es una
sesión de apertura, pero mostró cómo es el feed gratuito:
\begin{itemize}
\item llegaron cotizaciones para los 3\,250 contratos pedidos;
\item solo el 9--15\,\% de los precios de SPXW cae en la grilla de ticks de la bolsa, cuando en un NBBO real
caerían todos;
\item entre el 20 y el 28\,\% de los pares de SPXW viola la banda de la paridad;
\item en el modo descriptivo del cierre, que no es elegible para el piloto, el RR25 a 30 días fue de
$-3.08$ puntos en SPXW y de $-3.20$ en SPY, con una sonrisa coherente.
\end{itemize}
Con este feed, los residuos de paridad no sirven como señal ni como control fino de calidad, y sin OPRA no se
puede medir cuánto sesga el RR25. El piloto mide el feed \texttt{indicative}; la decisión sobre OPRA se toma
con las primeras sesiones reales.

\subsection{La captura diaria}\label{sec:estado_captura}

La captura corre en GitHub Actions, dentro de un repositorio \textbf{privado}: los datos de Alpaca y de OPRA no
se pueden redistribuir, y el repositorio del código es público. Los workflows descargan el código del
repositorio público en un commit fijo, hoy \texttt{4ffde8a}, y cada registro guarda el commit usado. Cambiar
ese commit es cambiar la versión de medición.

\begin{figure}[H]
\centering
\resizebox{\linewidth}{!}{% Jornada de captura en el repositorio de datos (hora de Nueva York, EDT). t = minutos desde las 09:00.
\begin{tikzpicture}[x=1mm,y=1cm,
  disparo/.style={circle,fill=#1,minimum size=2.5mm,inner sep=0pt},
  respaldo/.style={circle,draw=#1,fill=white,line width=0.8pt,minimum size=2.5mm,inner sep=0pt},
  ventana/.style={fill=#1!6,draw=#1!30,line width=0.4pt},
  espera/.style={fill=#1!28,draw=none},
  carril/.style={titulo,anchor=east,align=right}]
% Ventanas en que la compuerta deja arrancar
\filldraw[ventana=qazul] (0,2.42) rectangle (43,2.78);
\filldraw[ventana=qazuloscuro] (0,1.52) rectangle (58,1.88);
\filldraw[ventana=qnaranja] (76,0.62) rectangle (135,0.98);
% Esperas hasta el corte
\fill[espera=qazul] (11,2.52) rectangle (45,2.68);
\fill[espera=qazuloscuro] (21,1.62) rectangle (60,1.78);
% Apertura y cortes
\draw[tenue,dashed,line width=0.6pt] (30,0) -- (30,3.1);
\node[nota,anchor=south] at (30,3.1) {apertura};
\draw[qazul,line width=0.9pt] (45,0) -- (45,3.1);
\node[nota,anchor=south,text=qazul] at (45,3.1) {\textbf{corte}};
\draw[qazuloscuro,line width=0.9pt] (60,0) -- (60,3.1);
\node[nota,anchor=south,text=qazuloscuro] at (60,3.1) {\textbf{corte}};
% Disparos y respaldos
\node[disparo=qazul] at (11,2.6) {};
\node[respaldo=qazul] at (26,2.6) {};
\node[disparo=qazuloscuro] at (21,1.7) {};
\node[respaldo=qazuloscuro] at (36,1.7) {};
\node[disparo=qnaranja] at (81,0.8) {};
% Carriles
\node[carril] at (-2,2.6) {captura-0945};
\node[carril] at (-2,1.7) {captura-1000};
\node[carril] at (-2,0.8) {historico-alpaca};
% Notas de cada carril
\node[nota,anchor=west,align=left] at (64,2.6) {ráfaga en paralelo 5 s antes del corte de las 09:45;\\cada página se guarda al llegar};
\node[nota,anchor=west,align=left] at (64,1.7) {la misma lógica para el corte de las 10:00,\\en otra máquina y otro proceso};
\node[nota,anchor=north west,align=left] at (76,0.56) {SIP de SPY en cada corte (consultable 15 min después) y dividendos};
% Eje
\draw[guia,line width=0.7pt] (0,0) -- (150,0);
\foreach \t/\h in {0/09:00,15/09:15,30/09:30,45/09:45,60/10:00,75/10:15,90/10:30,105/10:45,120/11:00,135/11:15,150/11:30}{
  \draw[guia] (\t,0) -- (\t,-0.1);
  \node[nota,anchor=north] at (\t,-0.12) {\h};}
% Leyenda
\begin{scope}[yshift=-0.85cm]
\node[disparo=tinta2] at (2,0) {};
\node[nota,anchor=west] at (4,0) {disparo programado};
\node[respaldo=tinta2] at (32,0) {};
\node[nota,anchor=west] at (34,0) {respaldo, 15 min después};
\fill[espera=tinta2] (68,-0.08) rectangle (76,0.08);
\node[nota,anchor=west] at (77,0) {espera hasta el corte};
\filldraw[ventana=tinta2] (104,-0.18) rectangle (112,0.18);
\node[nota,anchor=west] at (113,0) {la compuerta deja arrancar};
\end{scope}
\end{tikzpicture}
}
\caption{Jornada de captura en hora de Nueva York (EDT). Cada hora de corte tiene su disparo y su respaldo,
en máquinas separadas, y una compuerta deja arrancar solo a los que llegan a tiempo. Los horarios son los
programados: GitHub puede retrasarlos, y los registros miden el arranque real.}
\label{fig:jornada}
\end{figure}

Cada corte termina en un estado: \texttt{completa}, \texttt{parcial}, \texttt{fallida} o, si no hubo captura,
\texttt{perdida}. Cada ejecución deja un registro, también cuando falla al iniciar o después de capturar.
Tres mecanismos protegen la captura:
\begin{itemize}
\item un diario por página permite recuperar una captura interrumpida;
\item un plazo de preparación deja llegar a tiempo al respaldo;
\item un plazo absoluto impide que un proceso colgado siga corriendo.
\end{itemize}
Solo se guarda el crudo: las tablas se reconstruyen a partir de él, y reprocesarlo da los mismos bytes. Las
credenciales nunca se escriben en un archivo.

\subsection{Objetivo, etiquetas y dividendos}\label{sec:estado_objetivo}

El piloto separa tres papeles que no se mezclan:
\begin{itemize}
\item \textbf{señal}: las opciones SPXW;
\item \textbf{referencia de las opciones}: el nivel implícito de SPX, que usan los controles y las medidas de
la cadena;
\item \textbf{objetivo}: SPY observado en el SIP, con el rendimiento total (con dividendos) y el de precio por
separado.
\end{itemize}

Las etiquetas miden 1 y 5 sesiones, con la misma hora inicial y final. Alpaca no garantiza cuándo publica un
dividendo, así que cada dividendo se guarda en versiones, con la fecha desde la que se conoce cada una. Si
una consulta ya no trae un dividendo, o lo trae sin fecha ex o sin monto, se abre una discrepancia. Las
etiquetas que podría afectar quedan \emph{pendientes} hasta que la resuelva evidencia fechada. Una etiqueta
se acepta cuando una consulta hecha al menos 60 días después de su fin no deja discrepancias abiertas: es
una regla, no una garantía. El entrenamiento solo usa las etiquetas que ya habrían madurado en la fecha
simulada.

\subsection{Revisiones externas}\label{sec:estado_revisiones}

Cada entrega pasó por una revisión externa. Cada hallazgo se reprodujo, se convirtió en una prueba de
regresión y se corrigió (\tabref{tab:revisiones}). Entre tanto, la batería pasó de 142 a 203 pruebas.

\begin{table}[H]
\centering\small
\caption{Rondas de revisión y su resultado.}
\label{tab:revisiones}
\begin{tabularx}{\linewidth}{@{}>{\RaggedRight}p{3.1cm} Y >{\RaggedRight\arraybackslash}p{2.5cm}@{}}
\toprule
\encab{Revisión} & \encab{Qué encontró o pidió} & \encab{Corrección}\\
\midrule
\texttt{e2b92f0} & Ocho hallazgos, entre ellos la elegibilidad estricta, los estados de captura y la
separación de fuentes. & \texttt{79c1554}\\
\texttt{d815bdd} & Separar señal, referencia y objetivo; objetivo SPY observado; dividendos; recuperación de
capturas. & \texttt{03f2230} y siguientes\\
\texttt{8c97b2b} & Cobertura de dividendos aparte; versiones de eventos y etiquetas; plazo para el respaldo. &
\texttt{6f0e748}, \texttt{f2a5663}\\
\texttt{a5e2748} & Una ausencia no retira un dividendo, sino que abre una discrepancia; estado puntual de las
etiquetas. & \texttt{c5aaded}, \texttt{8316540}\\
\texttt{5c15028} & Un dividendo sin fecha ex o sin monto deja pendientes las etiquetas afectadas. &
\texttt{7fb0585}\\
Cierre de \texttt{b240dc2} & Sin defectos nuevos. Después: registro de cada ejecución a prueba de errores y
revisión de la operación. & \texttt{7a59d04}, \texttt{bbbdc6e}\\
Cambios operativos & Registro también al fallar al iniciar; ruta explícita de resoluciones; pruebas en
Windows. & \texttt{4ffde8a}\\
Verificación de \texttt{4ffde8a} & Sin defectos; recomienda fijarlo en los workflows. & \texttt{e1a7832}\\
Instalación & Verificación independiente del repositorio de datos y de la prueba manual. & Confirmada\\
\bottomrule
\end{tabularx}
\end{table}

\subsection{Instalación y aceptación}\label{sec:estado_instalacion}

El repositorio de datos se instaló la noche del 27 de septiembre, con las plantillas revisadas, idénticas
byte a byte a las del commit \texttt{e1a7832}. La prueba manual salió bien:
\begin{itemize}
\item la captura inmediata quedó \texttt{completa}, con 3\,250 cotizaciones, y su registro guarda el commit
\texttt{4ffde8a} sin cambios locales;
\item el histórico descargó completo el SIP de SPY en los 10 cortes del 21 al 25 de septiembre, además de los
dividendos;
\item no apareció ninguna credencial en los datos guardados.
\end{itemize}
Una verificación independiente descargó el repositorio y reprodujo los mismos resultados.

\begin{criterio}[Criterio para aceptar la operación]
Se revisan 3--5 sesiones reales, desde el 28 de septiembre, y en los dos cortes:
\begin{itemize}
\item el estado, el retraso del arranque, la preparación, el margen de la ráfaga y las respuestas tardías;
\item el comportamiento del respaldo cuando actúe;
\item el SIP y los dividendos, no solo el código de salida;
\item que las capturas del horario regular sean elegibles para el piloto y que cada registro conserve el
commit fijado.
\end{itemize}
\end{criterio}

\begin{notatecnica}[Mantenimiento antes del 19 de octubre]
GitHub pasará la etiqueta \texttt{ubuntu-latest} de Ubuntu 24.04 a 26.04 de forma gradual entre el 19 de
octubre y el 19 de noviembre \citep{github2026}. Eso cae en medio del piloto, y unas corridas usarían una
imagen y otras, la otra. Antes de esa fecha, las plantillas se fijarán en \texttt{ubuntu-24.04}, la imagen que
ya se usa, junto con versiones de las acciones hechas para Node.js 24. El cambio se instalará con una corrida
manual de prueba.
\end{notatecnica}

\subsection{Qué falta}\label{sec:estado_falta}

La \figref{fig:ruta} ordena lo que queda. Cada paso exige que el anterior haya superado su criterio.

\begin{figure}[H]
\centering
\resizebox{\linewidth}{!}{% Estado y ruta al 28 de septiembre de 2026: qué está hecho, qué sigue y qué queda para después.
\begin{tikzpicture}[
  col/.style={cajallena=#1,text width=2.75cm,minimum height=9mm,font=\sffamily\scriptsize\bfseries},
  item/.style={caja=#1,text width=2.75cm,minimum height=7mm,font=\sffamily\tiny},
  node distance=1.6mm and 2.4mm]
\node[col=bueno] (c1) {HECHO};
\node[col=qazul,right=of c1] (c2) {AHORA};
\node[col=qvioleta,right=of c2] (c3) {PRÓXIMO};
\node[col=qnaranja,right=of c3] (c4) {DESPUÉS};
\node[col=tenue,right=of c4] (c5) {MÁS ADELANTE};
% Hecho
\node[item=bueno,below=of c1] (h1) {Propuesta y núcleo de referencia};
\node[item=bueno,below=of h1] (h2) {Pipeline del piloto con datos sintéticos};
\node[item=bueno,below=of h2] (h3) {Fuente de datos verificada: Alpaca};
\node[item=bueno,below=of h3] (h4) {Objetivo SPY, etiquetas y dividendos};
\node[item=bueno,below=of h4] (h5) {Captura diaria instalada y verificada};
% Ahora
\node[item=qazul,below=of c2] (a1) {Aceptación: 3--5 sesiones, del 28 de septiembre al 2 de octubre};
\node[item=qazul,below=of a1] (a2) {Mantenimiento de las plantillas antes del 19 de octubre};
\node[item=qazul,below=of a2] (a3) {Decisión sobre OPRA};
% Próximo
\node[item=qvioleta,below=of c3] (p1) {Piloto de 20--40 sesiones (hasta el 23 de octubre o el 20 de noviembre)};
\node[item=qvioleta,below=of p1] (p2) {Informe y dictamen de los datos};
\node[item=qvioleta,below=of p2] (p3) {Umbrales congelados en una versión nueva};
% Después
\node[item=qnaranja,below=of c4] (d1) {Protocolo predictivo fijado antes de mirar};
\node[item=qnaranja,below=of d1] (d2) {Historia ampliada: más de un año o comprada};
\node[item=qnaranja,below=of d2] (d3) {Prueba predictiva con periodo reservado};
% Más adelante
\node[item=tenue,below=of c5] (m1) {Cuantiles y transporte (variante F)};
\node[item=tenue,below=of m1] (m2) {Confianza, conformal y régimen};
\node[item=tenue,below=of m2] (m3) {Decisión sobre la posición};
\node[nota,text width=2.75cm,below=2mm of m3] {cada bloque entra solo si el anterior superó su criterio};
\end{tikzpicture}
}
\caption{Estado y ruta al 28 de septiembre de 2026.}
\label{fig:ruta}
\end{figure}

\begin{enumerate}
\item \textbf{Aceptación de la captura}, del 28 de septiembre al 2 de octubre: revisar 3--5 sesiones con el
criterio anterior.
\item \textbf{Mantenimiento de las plantillas}, antes del 19 de octubre.
\item \textbf{Piloto de medición}, de 20--40 sesiones: la sesión 20 es el 23 de octubre y la 40, el 20 de
noviembre. Su salida es un informe reproducible con un dictamen de los datos: «apto para ampliar historia»,
«apto con limitaciones» o «insuficiente». Cualquier relación con los rendimientos que se vea ahí es
exploratoria.
\item \textbf{Protocolo predictivo}, escrito antes de mirar el periodo reservado
(\secref{sec:estado_cambios}).
\item \textbf{Historia ampliada.} Entre 20 y 40 sesiones no bastan para contrastar la hipótesis. Si el efecto
es pequeño, hará falta más de un año de datos, salvo que se compre historia (por ejemplo, a Cboe DataShop).
\item \textbf{Después}, y solo con evidencia: cuantiles y transporte, confianza y la decisión sobre la
posición.
\end{enumerate}

\begin{pendiente}
\begin{itemize}
\item Contratar OPRA (99~USD al mes) o seguir con el feed gratuito, según las primeras sesiones.
\item Alargar el piloto de 20 a 40 sesiones.
\item Comprar historia para la prueba predictiva o acumularla hacia adelante.
\end{itemize}
\end{pendiente}

\subsection{Qué cambió respecto de esta propuesta}\label{sec:estado_cambios}

\begin{itemize}
\item \textbf{Fuente de datos.} La \tabref{tab:fuentes} dejaba a Alpaca para explorar; hoy es la fuente del
piloto, con captura hacia adelante porque no guarda cotizaciones pasadas. Cboe DataShop sigue como
candidata para la historia ampliada.
\item \textbf{Objetivo.} Alpaca no publica el nivel de SPX, así que el objetivo es SPY observado en el SIP, con
rendimiento total. El SPX implícito queda solo como referencia de las opciones.
\item \textbf{Primera pregunta.} La hipótesis general de la \secref{sec:hipotesis} se concretó para la primera
prueba. En la ecuación~\eqref{eq:diferencial}, la señal es el cambio del RR25 respecto de la sesión anterior,
el objetivo es el signo del rendimiento de SPY a una sesión y la puntuación es la de Brier. El horizonte de
cinco sesiones y el riesgo de cola quedan como análisis secundarios.
\item \textbf{Comparación por pasos.} Con la misma familia de modelos, se agrega una fuente de información
cada vez: primero la base (rendimientos pasados, salto de apertura, volatilidad reciente y calendario);
luego la volatilidad implícita, el nivel y el cambio de la asimetría y la paridad; y, al final, los
cuantiles y el transporte.
\item \textbf{Paridad.} Con el feed gratuito, los residuos de paridad son un diagnóstico de calidad, no una
señal.
\end{itemize}

\begin{noconcluirfinal}
\begin{itemize}
\item La instalación prueba que la captura funciona, no que la señal sirva.
\item La primera verificación con datos reales es del cierre de un viernes, no de una apertura.
\item El piloto evaluará la calidad de los datos, no la capacidad predictiva. Todavía no hay base para
recomendar posiciones.
\end{itemize}
\end{noconcluirfinal}
